%-----Important Packages-----
\usepackage{setspace}
\usepackage[table, dvipsnames]{xcolor}
\usepackage{mathtools}
\usepackage{amssymb}
\usepackage{amsthm}
\usepackage{gensymb}
\usepackage{wasysym}
\usepackage{longtable}
\usepackage{array}
\usepackage{caption}
\usepackage{circuitikz}
\usepackage{datetime}
\usepackage{enumitem}
\usepackage{etoolbox}
\usepackage{fancyhdr}
\usepackage{float}
\usepackage{fontspec}
\usepackage{geometry}
\usepackage{graphicx}
\usepackage[hidelinks]{hyperref}
\usepackage{listings}
\usepackage{pifont}
\usepackage{physics}
\usepackage{ragged2e}  % better centering
\usepackage{soul}
\usepackage[most]{tcolorbox}
\usepackage{tikz}
\usepackage{titlesec}
\usepackage{tocloft}
\usepackage{tabularx}
\usepackage{xparse}
\usepackage[scaled=0.95]{helvet}
\usepackage{booktabs}
\usepackage{upquote}
\usepackage{ltablex}
\usepackage[nameinlink]{cleveref}
\usepackage{nameref}
\usepackage{xurl}
\usepackage{seqsplit}
% \usepackage{minted}



\tcbuselibrary{skins}

\newtcbox{\codechip}{
  colback=gray!15, colframe=gray!15,
  arc=2pt, boxsep=0pt,
  left=3pt, right=3pt, top=1pt, bottom=1pt,
  fontupper=\ttfamily
}
\newcommand{\seqtt}[1]{\texttt{\seqsplit{#1}}}


\usetikzlibrary{shapes.geometric, arrows.meta, positioning, calc}
\keepXColumns

% Creating a brand new command named \normalpath
\DeclareUrlCommand\normalpath{\small \urlstyle{same}}

\hypersetup{
    colorlinks=true,
    linkcolor=blue!50!black,
    citecolor=blue!50!black,
    urlcolor=blue!60!black
}

%-----Page Layout and Spacing-----
% Configures page margins
\geometry{top=1in, bottom=.7in, left=.7in, right=.7in}


% Level 1: Standard Solid Bullet
\renewcommand{\labelitemi}{$\bullet$}
% Level 2: Your suggested Pointer (⮚ style)
\renewcommand{\labelitemii}{\ding{226}}
% Level 3: Small Hollow Circle
\renewcommand{\labelitemiii}{$\circ$}

% Custom column types for tables
\renewcommand{\arraystretch}{1.3}
\AddToHook{env/tabularx/begin}{\rowcolors{2}{gray!15}{white}}
\setlength{\arrayrulewidth}{0.8pt}
% 1. Vertically center all 'X' columns (Middle alignment):
\renewcommand{\tabularxcolumn}[1]{m{#1}}
% 2. Define 3 custom column types for horizontal alignment:For fixed-width columns with specific horizontal alignment:
\newcolumntype{x}[1]{>{\raggedright\arraybackslash}m{#1}}
\newcolumntype{y}[1]{>{\centering\arraybackslash}m{#1}}
\newcolumntype{z}[1]{>{\raggedleft\arraybackslash}m{#1}}
% 2. Define 3 custom column types for horizontal alignment: for flexible-width columns that expand to fill the available space:
\newcolumntype{C}{>{\centering\arraybackslash}X}    % C = Centered
\newcolumntype{L}{>{\raggedright\arraybackslash}X}  % L = Left-aligned
\newcolumntype{R}{>{\raggedleft\arraybackslash}X}   % R = Right-aligned
% Individual cell overriding commands for Headers
% For flexible X-columns (Centered)
\newcommand{\CC}[1]{\multicolumn{1}{|C|}{\textbf{#1}}}
% For fixed-width columns (Centered)
\newcommand{\yy}[2]{\multicolumn{1}{|y{#2}|}{\textbf{#1}}}
% Standard content-width (Centered) 
\newcommand{\cc}[1]{\multicolumn{1}{|c|}{\textbf{#1}}}


% Emoji/symbol font: Segoe UI Emoji on Windows; DejaVu Sans elsewhere (Linux/Claude builds)
\IfFontExistsTF{Segoe UI Emoji}{\newfontfamily\emojifont{Segoe UI Emoji}}{%
  \IfFontExistsTF{DejaVu Sans}{\newfontfamily\emojifont{DejaVu Sans}}{\newcommand{\emojifont}{}}}

\graphicspath{{figures/}}
% Custom Counter for question block
\newcounter{shadowq}
% \setstretch{0.7}

\let\oldtextbf\textbf
\renewcommand{\textbf}[1]{{\color{green!55!black} \oldtextbf{#1}}}

%Format Spacing. Command: \titlespacing*{<sectioning command>}{<left>}{<before-sep>}{<after-sep>}
\titlespacing*{\section}{0pt}{15pt}{-5pt}
\titlespacing*{\subsection}{0pt}{10pt}{-5pt}
\titlespacing*{\subsubsection}{0pt}{5pt}{-4pt}


% Paragraph spacing
\setlength{\parskip}{0.5em}   % space between paragraphs
\setlength{\parindent}{0pt}   % no indent
% \setlength{\intextsep}{-5pt}  % space above and below in-line figures
\setlength{\intextsep}{-3pt plus 3pt minus 2pt}   % when table remain between text [h]
\setlength{\textfloatsep}{0pt  plus 3pt minus 2pt} % when table remain at the top [t] or bottom [b] of the page
\setlength{\floatsep}{0pt plus 3pt minus 2pt}     % gap between two consecutive tables on the same page

% Customizes caption formatting for figures and tables
\captionsetup{margin=1cm, font=small, labelfont=bf, justification=centering,}
\setlength{\abovecaptionskip}{5pt} % Removes the unwanted whitespace after the caption
\setlength{\belowcaptionskip}{5pt}  % Adds space between table bottom and caption text

% Bullet list spacing
\setlist[itemize]{itemsep=0em, topsep=.2em}
\setlist[enumerate]{itemsep=0em, topsep=.2em}

% Customizes the header and footer layout for a professional look.
\pagestyle{fancy}
\fancyhf{}
% Header text is set per document in main.tex with \newcommand{\ReportHeader}{...}
\providecommand{\ReportHeader}{}
\fancyhead[L]{\ReportHeader}
\fancyhead[R]{\hyperlink{toc}{\fbox{\small TOC}}}
\fancyfoot[C]{\thepage}
\setlength{\headheight}{0pt}

% ---------- Question box ----------
\newtcolorbox{mainquestion}{
	before=\clearpage,
	colback=gray!60,        % light gray background
	colframe=black,       % subtle border
	coltext=black,         % black text
	boxrule=0.6pt,
	arc=4pt,
	left=8pt,
	right=8pt,
	top=6pt,
	bottom=6pt
}

\definecolor{codegreen}{HTML}{228B22}
\definecolor{codepurple}{HTML}{8F00FF}
\definecolor{codered}{HTML}{B22222}
\definecolor{codegray}{HTML}{F5F5F5}
\definecolor{linegray}{HTML}{999999}
\definecolor{framegray}{HTML}{CCCCCC}
\colorlet{numb}{yellow!80!green!80!black}
\definecolor{BlueLagoon}{HTML}{007F8C}
\definecolor{DeepOcean}{HTML}{005F73}

\lstset{
	basicstyle=\ttfamily\small,              % Base font
	keywordstyle=\color{codepurple}\bfseries, % Keywords
	commentstyle=\color{codegreen}\itshape, % Comments
	stringstyle=\color{codered},            % Strings
	backgroundcolor=\color{codegray},       % Background
	numbers=left,                           % Left-side numbers
	numberstyle=\ttfamily\small\bfseries\color{linegray},      % Number style
	stepnumber=1,                           % Every line numbered
	numbersep=10pt,                         % Gap from code
	frame=tbl,                               % Top & bottom lines
	framerule=0.8pt,                        % Frame thickness
	rulecolor=\color{framegray},            % Frame color
	framesep=3pt,                           % Inner padding
	breaklines=true,                        % Auto line wrap
	keepspaces=true,						% Preserve spaces
	breakatwhitespace=true,                 % Break at spaces
	tabsize=4,                              % Tab width
	columns=fullflexible,
	extendedchars=true,
	upquote=true,						    % Straight quotes
	showspaces=false,                       % Hide space markers
	showtabs=false,                         % Hide tab markers
	showstringspaces=false,                 % Hide string spaces
	xleftmargin=25pt,                       % Left margin
	xrightmargin=0pt,                       % Right margin
	framexleftmargin=0pt,					% Frame left margin
	aboveskip=15pt,                         % Space above
	belowskip=15pt,                         % Space below
	captionpos=b,                           % Caption bottom
	% Wrapped line marker
	postbreak=\mbox{\textcolor{linegray}{\ensuremath{\hookrightarrow}\space}},
	escapeinside={(*@}{@*)},                % Inline LaTeX escape, use like: (*@$1, 2, 3, 4, \dots$@*)
	literate={~}{{$\sim$}}1,
}

% --- Step 2: Defining Language-Specific Styles ---
% Style for C++
\lstdefinestyle{myCPP}{
	language=C++,
	keywordstyle=\color{codepurple}\bfseries,
	commentstyle=\color{codegreen}\itshape,
	stringstyle=\color{codered},
}

\lstdefinelanguage{json}{
literate=
*{0}{{{\color{numb}0}}}{1}
{1}{{{\color{numb}1}}}{1}
{2}{{{\color{numb}2}}}{1}
{3}{{{\color{numb}3}}}{1}
{4}{{{\color{numb}4}}}{1}
{5}{{{\color{numb}5}}}{1}
{6}{{{\color{numb}6}}}{1}
{7}{{{\color{numb}7}}}{1}
{8}{{{\color{numb}8}}}{1}
{9}{{{\color{numb}9}}}{1}
{:}{{{\color{red}{:}}}}{1}
{,}{{{\color{red}{,}}}}{1}
{\{}{{{\color{blue}{\{}}}}{1}
{\}}{{{\color{blue}{\}}}}}{1}
{[}{{{\color{blue}{[}}}}{1}
{]}{{{\color{blue}{]}}}}{1},
}
% % Minted settings (if you choose to use minted instead of listings)
% \renewcommand{\theFancyVerbLine}{\textcolor{linegray}{\footnotesize\ttfamily\arabic{FancyVerbLine}}}
% \usemintedstyle{friendly}
% \setminted{
% 	style=default,
% 	fontsize=\small,
% 	linenos=true,
% 	numbers=left, 
% 	stepnumber=1,
% 	numbersep=5pt,
% 	bgcolor=codegray,
% 	frame=single,
% 	framerule=0.8pt,
% 	rulecolor=framegray,
% 	framesep=3pt,
% 	breaklines=true,
% 	breakbytoken=true,
% 	tabsize=4,
% 	xleftmargin=12pt,
% 	xrightmargin=0pt,
% 	stripnl=true,
% 	breakbytoken=true,
% }


% ---------- Highlight commands ----------
\newcommand{\hlblack}[1]{%
	\tcbox[
		on line, % Will stay within the line of text
		arc=2pt, % How rounded the corners will be (you can change this)
		colback=black!95, % Background color
		colframe=black!95, % Border color
		boxsep=0pt, left=2pt, right=2pt, top=2pt, bottom=2pt
	]{\textcolor{white}{\textbf{#1}}}%
}

\newcommand{\grayhl}[1]{{\sethlcolor{gray!50}\color{black}{\hl{#1}}}}
\newcommand{\blackhl}[1]{{\sethlcolor{black!80}\color{white}{\hl{#1}}}}
\newcommand{\note}[1]{\emph{\scriptsize (#1)}}
\newcommand{\keyterm}[1]{{\normalfont\bfseries #1}}


% Header Formatting
\newcommand{\addcolon}[1]{#1:}
% ---------- Section format ----------
\titleformat{\section}
{\normalfont\Large\bfseries\color{blue!80!black}}
{\thesection.}
{0.3em}
{\addcolon}
% ---------- Subsection format ----------
\titleformat{\subsection}
{\normalfont\large\bfseries\color{red!75!black}}
{\thesubsection.}
{0.3em}
{\addcolon}
% ---------- Subsubsection format ----------
\titleformat{\subsubsection}
{\normalfont\normalsize\bfseries\color{blue!80!black!80!white}}
{\thesubsubsection.}
{0.3em}
{\addcolon}

% Custom environment for displaying equations/content with an optional side note section
\newenvironment{eqnote}[2]
{
	\begingroup
	\par
	\setlength{\parskip}{4pt}
	\setlength{\abovedisplayskip}{0.0em}
	\setlength{\belowdisplayskip}{0.0em}
	\noindent
	\ifstrempty{#2}
	{%
		% --- If note is empty ---
		#1
	}
	{%
		% --- If note exists ---
		\begin{minipage}[c]{0.65\textwidth}
			#1
		\end{minipage}%
		\hfill
		\begin{minipage}[c]{0.32\textwidth}
			{\scriptsize
			\vrule width 0.5pt
			\hspace{0.0em} % Slightly increased spacing for better readability
			\parbox[t]{.95\linewidth}{
			\textbf{Note,}\\[0.1em]
			\begingroup
			\setlength{\parindent}{0pt}
			\renewcommand{\baselinestretch}{1}
			\def\\{\par\textbullet\hspace{0.0em}}
			\textbullet\ #2
			\endgroup
			}
			}
		\end{minipage}
		\vspace{3.5pt}
	}
}
{
	\par
	\endgroup
}


\newtcolorbox{repocard}[1]{
    enhanced,
    colback=gray!15,
    colframe=black!60,
    coltitle=white,
    colbacktitle=BlueLagoon,
    fonttitle=\bfseries,
    title=#1,
    arc=2mm
}

% =====================================================================
% ---------- Study-notes additions (callouts, links, PDF navigation) ----------
% Added for exam notes. Everything above this line is the original house preamble.
% =====================================================================

% PDF navigation: bookmarks panel open on launch, numbered.
% (main.tex passes hyperref's destlabel option so another PDF can link here with \href{file.pdf\#label}{...}.)
\hypersetup{
	bookmarksopen=true,
	bookmarksopenlevel=1,
	bookmarksnumbered=true,
	pdfpagemode=UseOutlines
}

\colorlet{calloutnote}{black!65}
\colorlet{calloutexample}{DeepOcean}
\colorlet{calloutexam}{orange!80!black}
\colorlet{calloutwarn}{red!65!black}
\colorlet{calloutquestion}{BlueLagoon}
\colorlet{calloutsolution}{green!40!black}

% Shared look: same family as repocard (rounded, coloured title bar), breakable across pages.
\tcbset{
	calloutbase/.style={
		enhanced, breakable,
		arc=2mm, boxrule=0.6pt,
		left=6pt, right=6pt, top=4pt, bottom=4pt,
		toptitle=1pt, bottomtitle=1pt,
		fonttitle=\bfseries\small, coltitle=white,
		before skip=8pt, after skip=8pt,
		parbox=false
	}
}

% Note / notice (📌, plain blockquotes)
\newtcolorbox{notebox}[1][Note]{calloutbase, colback=gray!8, colframe=calloutnote, colbacktitle=calloutnote, title={#1}}
% Worked example
\newtcolorbox{examplebox}[1][Example]{calloutbase, colback=DeepOcean!5, colframe=calloutexample, colbacktitle=calloutexample, title={#1}}
% Exam focus (🎯)
\newtcolorbox{exambox}[1][Exam Focus]{calloutbase, colback=orange!7, colframe=calloutexam, colbacktitle=calloutexam, title={#1}}
% Mistake / correction / caution (⚠)
\newtcolorbox{warnbox}[1][Watch Out]{calloutbase, colback=red!4, colframe=calloutwarn, colbacktitle=calloutwarn, title={#1}}
% Practice question: {title}{label of its solution}; title bar carries a "Solution" link
\newtcolorbox{questionbox}[2]{calloutbase, colback=white, colframe=calloutquestion, colbacktitle=calloutquestion,
	title={#1\hfill\ifstrempty{#2}{}{\hyperref[#2]{\color{white}\normalfont\small Solution $\rightarrow$}}}}
% Solution: {title}{label of its question}; title bar carries a "Back" link
\newtcolorbox{solutionbox}[2]{calloutbase, colback=green!3, colframe=calloutsolution, colbacktitle=calloutsolution,
	title={#1\hfill\ifstrempty{#2}{}{\hyperref[#2]{\color{white}\normalfont\small $\leftarrow$ Back to question}}}}
% Self-check list at the end of a lecture
\newtcolorbox{selfcheckbox}[1][Self-check]{calloutbase, colback=BlueLagoon!4, colframe=calloutquestion, colbacktitle=calloutquestion, title={#1}}

% Inline markers that replace emoji the main font cannot draw
\newcommand{\exammark}{\tcbox[on line, arc=1.5pt, colback=calloutexam, colframe=calloutexam, boxsep=0pt,
	left=2pt, right=2pt, top=1pt, bottom=1pt]{\color{white}\bfseries\scriptsize EXAM}}
\newcommand{\pinmark}{\tcbox[on line, arc=1.5pt, colback=calloutnote, colframe=calloutnote, boxsep=0pt,
	left=2pt, right=2pt, top=1pt, bottom=1pt]{\color{white}\bfseries\scriptsize NOTE}}
\newcommand{\warnmark}{{\color{calloutwarn}\emojifont ⚠}}
\newcommand{\okmark}{{\color{green!45!black}\emojifont ✔}}
\newcommand{\nomark}{{\color{red!70!black}\emojifont ✘}}

% Small clickable chip for "go to" links: \golink{label}{text}
\newcommand{\golink}[2]{\hyperref[#1]{\tcbox[on line, arc=1.5pt, colback=BlueLagoon!12, colframe=BlueLagoon!60,
	boxsep=0pt, left=3pt, right=3pt, top=1pt, bottom=1pt]{\small\color{BlueLagoon!70!black}#2}}}

% Figures: the house \intextsep is negative (keeps [h] tables tight), which pulls [H] figures
% up into the preceding text/heading. Give figures their own positive spacing; tables unchanged.
\AddToHook{env/figure/begin}{\setlength{\intextsep}{10pt plus 2pt minus 2pt}}

% Table cells holding tall maths (fractions, integrals, matrices): \tallcell{...} pads the cell
% 3pt above and below so the formula never touches the row rules. Only for x/y/z and X-type
% (paragraph) columns; the converter never puts tall maths in plain l/c/r columns.
\newcommand{\tallcell}[1]{\vspace*{3pt}#1\par\vspace*{3pt}}

% Keep a heading on the same page as the table (or figure) right after it.
% The converter writes \needspace{...} before such headings.
\usepackage{needspace}
